% Blank IEEE conference manuscript based on conference_101719.tex.
% Keep this file alongside IEEEtran.cls. Replace the placeholders and add your text.
\documentclass[conference]{IEEEtran}

% Packages retained from the supplied template.
\usepackage{cite}
\usepackage{amsmath,amssymb,amsfonts}
\usepackage{algorithmic}
\usepackage{graphicx}
\usepackage{textcomp}
\usepackage{xcolor}
\usepackage{xurl}

% Allow a tall two-column table to occupy more of the top of a page.
\renewcommand{\dbltopfraction}{0.9}
\renewcommand{\floatpagefraction}{0.75}

\begin{document}

\title{Distributed Optimization for Day-Ahead Scheduling in Local Energy Systems: A Review}

\author{\IEEEauthorblockN{Yaomin Tang}
\IEEEauthorblockA{%\textit{Your Department} \\
\textit{RWTH Aachen University} \\
Aachen, Germany \\
yaomin.tang@rwth-aachen.de}}

\maketitle

\begin{abstract}
% Write the abstract here.
As rooftop photovoltaic (PV) systems and flexible loads become more common among customers connected to distribution networks, some customers are becoming prosumers who both generate and consume electricity, while flexible loads allow users to adjust their demand. Variations in renewable generation and demand complicate day-ahead scheduling in local energy systems, where aggregators must coordinate flexible resources across multiple participants. This paper reviews centralized and distributed optimization approaches to this scheduling problem, focusing on dual decomposition (DD) and the alternating direction method of multipliers (ADMM). Drawing on existing studies, it examines how these methods are applied to portfolio balancing and compares their reported characteristics in terms of solution quality, computational performance, scalability, and information exchange. The review discusses the strengths and limitations of distributed optimization relative to centralized approaches.
 

\end{abstract}

\begin{IEEEkeywords}
% Add keywords separated by commas.
Local energy systems, day-ahead scheduling, operational flexibility, distributed optimization, distributed energy resources
\end{IEEEkeywords}

\section{Introduction}
% Write the introduction here.
Rooftop photovoltaic (PV) systems, residential batteries, heat pumps, and electric-vehicle chargers are behind-the-meter resources typically owned and operated by different customers \cite{Diekerhof2018,OlivellaRosell2020}. Metering at the grid connection reveals the net power exchange but not, by itself, the operating states, constraints, or preferences of the individual devices \cite{Diekerhof2018}, and uncertainty in PV output and demand further complicates the forecasting and maintenance of day-ahead schedules \cite{Diekerhof2018}. Scheduling these resources therefore requires coordinating customer-side flexibility under limited information sharing, while respecting both economic objectives and device constraints.

The difficulty lies in the coupling: the schedules chosen by individual customers must sum to the portfolio that the aggregator commits to, or to the flexibility volume that the distribution system operator requests. A single entity can in principle solve the coupled problem to optimality, but only if it holds the cost functions and device constraints of all participants \cite{Diekerhof2018}. Distributed algorithms relax this requirement: the computation is shared among many agents that obtain only the parameters they need through communication with a limited set of neighbours, which can also respect the privacy of data, cost functions, and constraints \cite{Molzahn2017}. Two decomposition schemes from this family, dual decomposition (DD) and the alternating direction method of multipliers (ADMM), have been applied specifically to day-ahead scheduling in local energy systems \cite{Diekerhof2018,OlivellaRosell2020}. The resulting evidence, however, is fragmented: reported computation times are obtained on different test systems and are therefore not directly comparable; the convergence guarantees reported are usually established for convex formulations, whereas the device models used in day-ahead scheduling often contain binary decisions such as the exclusivity of charging and discharging; such nonconvexities are addressed through convex relaxations, through exact local mixed-integer subproblems that forgo the convergence guarantee, or by applying the algorithms heuristically, and what each remedy costs in feasibility, optimality, or attainable flexibility is not reported in a consistent form. It is thus unclear how an aggregator should choose among these methods for a given scheduling problem, or whether a result reported for one setting carries over to another.

This review is restricted to decomposition methods that coordinate participants by exchanging their own schedules and a price for the shared constraint, while device-level schedules, local constraints, and cost functions remain local. Metaheuristic and learning-based approaches, such as particle swarm optimization and reinforcement learning, are not covered, nor are online or feedback-based dispatch schemes, in which set points are issued before the algorithm has converged; the computation of complete flexibility envelopes for aggregated resources is likewise outside its scope. Within this scope, the review makes three contributions. First, it presents the reviewed formulations in a common template and tabulates how they differ in objective structure, coupling constraint, and local feasibility set, so that their assumptions can be compared explicitly rather than inferred from separate case studies. Second, it assesses centralized optimization, DD, and ADMM in terms of effectiveness, scalability, and practical applicability to day-ahead scheduling in local energy systems, distinguishing them by the exchanged information and the guarantee they provide for reaching an optimum. Third, it derives conditional guidance for selecting an approach in a given setting and identifies where the available evidence is insufficient.

\section{Operational Flexibility and Day-Ahead Scheduling Problem}
% Define the local energy system and the day-ahead scheduling problem.

\subsection{Operational Flexibility and Enabling Conditions}
In this review, operational flexibility refers to the ability of a system to coordinate controllable assets across temporal and spatial scales so as to maintain a secure and economically efficient balance between supply and demand \cite{DallAnese2017}. Energy storage and demand-side management are widely regarded as the most promising sources of operational flexibility in power systems \cite{Diekerhof2018}. Within local energy systems, controllable assets mainly comprise battery and thermal energy storage, where the inherent thermal inertia of buildings can act as a passive thermal buffer, together with electric vehicles, heat pumps, and time-shiftable appliances. Unlocking this potential is nevertheless constrained by several physical and operational barriers: the limited availability of storage, particularly thermal storage; differences of several orders of magnitude between the time constants of electrical and thermal subsystems, which lead to severely mismatched response speeds; and the narrow set of flexible end-use options available in residential systems, which are largely confined to heating, ventilation, and air-conditioning equipment and schedulable household appliances \cite{DallAnese2017}. As summarized in \cite{Diekerhof2018}, operational flexibility can be characterized along four dimensions: power capacity, ramp rate, energy capacity, and duration.

Because residential flexibility is dispersed across numerous individual customers, coordinating it directly entails substantial data volumes and raises privacy concerns. Households and small-to-medium enterprises therefore commonly outsource their balancing responsibility to an aggregator, which schedules the pooled resources and offers aggregated flexibility services to distribution system operators (DSOs) and balance responsible parties (BRPs) \cite{Diekerhof2018,OlivellaRosell2020}. At the metering and settlement level, three conditions must be satisfied: certified metering data are needed to attribute energy flows to the correct participant and to settle them accordingly; financial adjustment and imbalance-responsibility rules must be defined between the aggregator and the BRP; and the flexibility service must be remunerated, for instance through tariff discounts or direct payments \cite{Diekerhof2018}. At the regulatory and control level, European electricity market unbundling prevents the objectives of the DSO and of end users from being merged and precludes sharing grid information with aggregators; as a result, the DSO typically quantifies its flexibility demand and the aggregator schedules its portfolio in response \cite{OlivellaRosell2020}. The traffic light concept for network operation further delineates the authority across three grid states: in the green state the aggregator schedules on market terms while the DSO merely monitors the network; in the yellow state the DSO intervenes actively and procures local flexibility through market-based incentives, compensating the aggregator for the portfolio capacity it forgoes; and in the red state, when security of supply is at risk, the DSO is legally entitled to control devices directly and may override schedules already issued by the aggregator \cite{Diekerhof2018}.

\subsection{Day-Ahead Optimal Dispatch and Portfolio Balancing Formulation}
\label{sec:formulation}

Let $\mathcal{N} = \{1,\dots,N\}$ denote the participants aggregated by one aggregator and $\mathcal{T} = \{1,\dots,T\}$ the periods of the day-ahead horizon. Participant $i$ decides a net power schedule $x_i \in \mathbb{R}^{T}$; the schedules admissible under its device dynamics, capacity limits, and user constraints form the set $\mathcal{X}_i$, and the cost incurred over the horizon is $f_i(x_i)$. This is the structure in which decomposition methods are usually stated \cite{Molzahn2017}, and the aggregator's problem reads
\begin{equation}
\begin{aligned}
\min_{x_1,\dots,x_N} \quad & \sum_{i \in \mathcal{N}} f_i(x_i) \\
\text{s.t.} \quad & x_i \in \mathcal{X}_i, && i \in \mathcal{N}, \\
& \sum_{i \in \mathcal{N}} A_i x_i = b .
\end{aligned}
\label{eq:template}
\end{equation}
For the aggregate balance considered here, $A_i$ selects the net power entries of $x_i$ and is the identity when $x_i$ is the net power schedule itself, whereas network constraints, where present, enter through non-identity $A_i$ and additional rows of $b$, provided that they are linear.

Whether \eqref{eq:template} is convex is decided by what is placed inside $\mathcal{X}_i$ and $f_i$, since the coupling constraint is linear by construction. Nonconvexity inside a local model may arise from discrete decisions such as the on/off state of a heat pump or the exclusivity of charging and discharging; from the underlying physics---power, gas, and water flows together with the heat-transfer models that couple them; or from the formulation itself, when continuous decision variables multiply one another \cite{DallAnese2017}. When the nonlinearity arises between participants instead, as in the ac power-flow equations relating neighbouring buses, convexity is lost, and much of the literature therefore works with linear approximations or convex relaxations of the flow equations; a coupling of that kind also falls outside the linear form of \eqref{eq:template} \cite{Molzahn2017}.

The template covers two settings that differ in where $b$ comes from. We distinguish them because the reviewed studies differ in exactly this respect. In autonomous scheduling, the aggregator optimises the portfolio schedule itself, so that $b$ is a decision variable \cite{Diekerhof2018}; in request-driven scheduling, $b$ is fixed by a flexibility request issued by the DSO or the BRP, and the aggregator must first establish whether the request can be met and at what cost \cite{OlivellaRosell2020}. Both settings share one structure: $f_i$ and $\mathcal{X}_i$ are held by participant $i$ and need not be disclosed, so that $b$ is the only quantity through which the participants are coupled \cite{Molzahn2017}.

Since $\mathcal{X}_i$ couples the periods within a day, the schedules that participant $i$ can deliver form a set $\mathcal{F}_i = \{A_i x_i : x_i \in \mathcal{X}_i\} \subset \mathbb{R}^{T}$ rather than an interval, and the portfolios that the participants can deliver jointly form the aggregate set $\mathcal{F} = \sum_{i \in \mathcal{N}} \mathcal{F}_i$, the Minkowski sum of the individual sets. A portfolio schedule $b$ is therefore attainable only if $b \in \mathcal{F}$. Describing $\mathcal{F}$ explicitly is impractical, as its dimension and the number of its facets grow with both $N$ and $T$; feasibility is therefore usually established by solving \eqref{eq:template} iteratively, which is where the methods reviewed next apply.


\section{Centralized and Distributed Optimization Approaches}

% Describe the selected approaches and the studies you compare.

\subsection{Centralized Optimization}

A centralized formulation solves \eqref{eq:template} as a single problem, in which one solver holds every $f_i$ and $\mathcal{X}_i$ and returns all $x_i$ together with $b$ \cite{Molzahn2017}. Because the device models contain binary decisions, such as the exclusivity of charging and discharging, the resulting problem is a mixed-integer program and is therefore nonconvex \cite{OlivellaRosell2020}; it is solved with a commercial branch-and-bound solver. Centralized computation nevertheless remains the natural reference, both as the baseline against which decomposed methods are judged and, in the traffic light concept, as the mode to which the DSO reverts in the red state \cite{Diekerhof2018}. How far it scales before the tree of binary decisions dominates the computation is taken up in Section~\ref{sec:discussion}.

\subsection{Dual Decomposition}

DD builds on the Lagrangian of \eqref{eq:template}. A multiplier $\lambda$ is introduced for the coupling constraint; it is common to all participants and admits the interpretation of a price for the coupled quantity. The resulting Lagrangian is

\begin{equation}
L(\{x_i\},\lambda) = \sum_{i\in\mathcal{N}} \big[ f_i(x_i) + \lambda^{\top} A_i x_i \big] - \lambda^{\top} b ,
\label{eq:dd-lag}
\end{equation}

Once the coupling constraint is absorbed into $L$, the remaining feasible set is the product $\mathcal{X} = \mathcal{X}_1 \times \cdots \times \mathcal{X}_N$, and every term of the sum depends on a single $x_i$ except $\lambda^{\top} b$, which depends on none of them. The minimum over a product set of such a sum equals the sum of the individual minima, so the dual function is

\begin{equation}
g(\lambda) = \sum_{i\in\mathcal{N}} \min_{x_i\in\mathcal{X}_i} \big[ f_i(x_i) + \lambda^{\top} A_i x_i \big] - \lambda^{\top} b ,
\label{eq:dd-dual}
\end{equation}

and each term can be evaluated by participant $i$ alone, using only its own objective, its own feasible set, and the common price $\lambda$. Since $g(\lambda)$ is a lower bound on the optimal value of \eqref{eq:template} for every $\lambda$, the dual problem is $\max_{\lambda} g(\lambda)$; $g$ is concave, being a pointwise minimum of affine functions, so this maximization is well posed. DD solves it iteratively: at iteration $k$, every participant solves its local problem at the current price $\lambda^{k}$, and the coordinator then adjusts the price in proportion to the residual,

\begin{equation}
x_i^{k+1} = \arg\min_{x_i\in\mathcal{X}_i} \big[ f_i(x_i) + (\lambda^{k})^{\top} A_i x_i \big], \quad \forall i\in\mathcal{N},
\label{eq:dd-local}
\end{equation}

\begin{equation}
\lambda^{k+1} = \lambda^{k} + \alpha^{k}\Big( \sum_{i\in\mathcal{N}} A_i x_i^{k+1} - b \Big).
\label{eq:dd-coord}
\end{equation}

Step \eqref{eq:dd-local} is a local scheduling problem: participant $i$ re-optimizes its own day-ahead schedule with an added linear term at the current price, and the $N$ problems are independent, so they can be solved in parallel. Step \eqref{eq:dd-coord} requires only the aggregate $\sum_i A_i x_i^{k+1}$, so the coordinator never sees a cost function or a device constraint. The update in \eqref{eq:dd-coord} then has a market reading: when the participants together deliver more than $b$, the price rises and their next local schedules shift accordingly \cite{Diekerhof2018}.

Optimality requires strong duality, which holds when the problem is convex; once binary decisions such as the exclusivity of charging and discharging are introduced, it does not hold in general \cite{Diekerhof2018}. Because the dual is generally nondifferentiable, \eqref{eq:dd-coord} is a subgradient step, and the iterates violate the coupling constraint until convergence, so only the converged schedule is used.

\subsection{Alternating Direction Method of Multipliers}

ADMM addresses the same template through a different splitting: it adds a copy of the local schedule and then minimizes over the copy and the original alternately. Participant $i$ keeps $x_i$ as its own variable and introduces $z_i$ as the copy that enters the coupling, so that

\begin{equation}
\min_{\{x_i\},\{z_i\}} \ \sum_{i\in\mathcal{N}} f_i(x_i) \quad \text{s.t.} \ x_i = z_i \ \forall i\in\mathcal{N}, \quad \sum_{i\in\mathcal{N}} A_i z_i = b ,
\label{eq:admm-form}
\end{equation}

which has the same solutions as \eqref{eq:template} whenever the latter is feasible. The copies do not change the problem; they create the second block that ADMM requires, because the penalty introduced below must measure a difference between two variables and \eqref{eq:template} offers only one. In place of the linear Lagrangian \eqref{eq:dd-lag}, the method uses an augmented Lagrangian \cite{Boyd2011} with one multiplier per copy constraint,

\begin{equation}
L_\rho = \sum_{i\in\mathcal{N}} \Big[ f_i(x_i) + \lambda_i^{\top}(x_i - z_i) + \frac{\rho}{2}\big\lVert x_i - z_i \big\rVert^{2} \Big],
\label{eq:admm-auglag}
\end{equation}

with a fixed penalty parameter $\rho>0$. The added penalty is what allows the two blocks to be minimized separately: at iteration $k$,

\begin{equation}
x_i^{k+1} = \arg\min_{x_i\in\mathcal{X}_i} \Big[ f_i(x_i) + (\lambda_i^{k})^{\top} x_i + \frac{\rho}{2}\big\lVert x_i - z_i^{k} \big\rVert^{2} \Big],
\label{eq:admm-x}
\end{equation}

and the copies follow by projecting the price-corrected schedules $x_i^{k+1} + \lambda_i^{k}/\rho$ onto the coupling constraint,

\begin{equation}
\begin{aligned}
\{z_i^{k+1}\} = {}&\arg\min_{\{z_i\}} \ \frac{\rho}{2}\sum_{i\in\mathcal{N}}\Big\lVert x_i^{k+1} + \frac{\lambda_i^{k}}{\rho} - z_i\Big\rVert^{2}\\
&\text{s.t.} \ \sum_{i\in\mathcal{N}} A_i z_i = b ,
\end{aligned}
\label{eq:admm-z-proj}
\end{equation}

a quadratic program whose cost grows with the number of scalar constraints---one per period in the day-ahead setting---rather than with the number of participants, because the copies enter only through the aggregate $\sum_i A_i z_i$. The multipliers then follow by

\begin{equation}
\lambda_i^{k+1} = \lambda_i^{k} + \rho\big( x_i^{k+1} - z_i^{k+1} \big).
\label{eq:admm-lambda}
\end{equation}

\begin{table*}[!t]
\centering
\caption{Coordination properties of the three approaches, compiled from \cite{Diekerhof2018,OlivellaRosell2020}. The optimum guarantee is a property of the method \emph{under a stated premise}, not of the method alone.}
\label{tab:comparison}
\footnotesize
\begin{tabular}{@{}p{0.17\textwidth}p{0.25\textwidth}p{0.25\textwidth}p{0.25\textwidth}@{}}
\hline
 & Centralized & DD & ADMM \\
\hline
Information reaching the coordinator
& cost functions and device constraints of all participants
& the schedules of all participants; returns one price for the coupling constraint
& the schedules of all participants; returns a corrected copy to each \\
Coordinator
& none, a single solver
& updates multipliers only
& projects the copies onto the shared constraint \\
Optimum guarantee
& optimum if the search terminates; the tree grows exponentially with the binary variables
& optimum if strong duality holds
& optimum for convex formulations; feasible, generally suboptimal, otherwise \\
Binary device decisions
& handled inside one solver
& convex relaxation; optimality then holds for the relaxed model only
& local mixed-integer solve, convex relaxation, or heuristic use; no convergence guarantee \\
\hline
\end{tabular}
\end{table*}

When every $A_i$ is the identity---that is, when the quantity exchanged is already the net exchange of participant $i$, with the device constraints absorbed in $\mathcal{X}_i$, as in all three studies reviewed here---the coupling constraint collapses to $\sum_{i\in\mathcal{N}} z_i = b$, and the projection collapses to subtracting the same offset from every price-corrected schedule; \eqref{eq:admm-z} states the result explicitly.

\begin{equation}
z_i^{k+1} = x_i^{k+1} + \frac{\lambda_i^{k}}{\rho} - \frac{1}{N}\Big( \sum_{j\in\mathcal{N}} x_j^{k+1} + \frac{1}{\rho}\sum_{j\in\mathcal{N}} \lambda_j^{k} - b \Big),
\label{eq:admm-z}
\end{equation}

so the coordinator computes one average and returns it: it solves no optimization problem and sees no cost function or device constraint, exactly as in \eqref{eq:dd-coord}. The local step \eqref{eq:admm-x} differs from \eqref{eq:dd-local} in carrying a proximal term that pulls the new schedule towards the previous copy, which is what keeps the iterates bounded when the price alone would not.

Because the copies are projected onto the coupling constraint at every iteration, each iterate satisfies $\sum_{i\in\mathcal{N}} A_i z_i = b$, so the residual that matters to the market vanishes from the start; the allocation nevertheless becomes a realizable device schedule only when $x_i$ and $z_i$ agree, which occurs at convergence. As with \eqref{eq:dd-coord}, the guarantee presupposes convex local models. A Jacobian variant instead has each participant replace the other participants' schedules by their values from the previous iteration, so that all participants can be solved concurrently; convergence then requires an extra proximal term that penalizes the distance to the participant's own previous schedule. This is the form used by Olivella-Rosell et al.\ \cite{OlivellaRosell2020}, and a fully distributed variant that updates \eqref{eq:admm-lambda} at the participants themselves is used by Lilla et al.\ \cite{Lilla2020}.


\section{Comparative Discussion}
\label{sec:discussion}

% DRAFT: the three criteria are defined here so that the comparison below does not
% rely on unstated notions of "better". Revise once Section III is written.

The criteria used below are fixed before the studies are compared. \emph{Effectiveness} denotes whether a feasible schedule of comparable quality is obtained: whether the returned plan satisfies the coupling constraint and the discrete device constraints, and how far its cost lies from the optimal cost of the centralized problem. \emph{Scalability} denotes how the largest problem that any single entity must solve, and the information that must be collected in one place, grow with the number of participants; for the iterative methods it also covers whether the number of iterations needed stays bounded. Wall-clock time is not used, since timings obtained on different test systems are not comparable. \emph{Practical applicability} denotes the conditions a method imposes on the setting in which it is to run: which data must be shared, whether a coordinator is available, and whether the device models are of a class for which the method can be used as intended, or only heuristically. Table~\ref{tab:comparison} collects the properties on which the three approaches differ.

% TABLE SKELETON. Cells deliberately state only what the cited studies support;
% leave a cell empty rather than filling it from general expectation.

The guarantee reported for a method depends on the model as much as on the algorithm, which the reviewed studies illustrate in opposite ways. Olivella-Rosell et al.\ solve each prosumer subproblem as a mixed-integer program, so that the exclusivity of charging and discharging is respected exactly, but the overall problem is nonconvex and the method is only expected to reach a feasible, generally suboptimal, portfolio \cite{OlivellaRosell2020}. Diekerhof instead replaces the nonconvex device models with convex approximations, which restores the convergence guarantee at the cost of model fidelity \cite{Diekerhof2018}. The two claims do not conflict: they hold under different premises, and comparing them by the label of the algorithm alone would be misleading. More decisive than the algorithm is therefore whether the discrete decisions are relaxed and whether the returned schedule is checked against them; the second of these is not reported in the reviewed studies, so the effectiveness of a method in a given setting cannot be inferred from the label attached to it.

The two methods also differ in what their guarantees demand of the parameters left to the user. In DD, the multiplier update \eqref{eq:dd-coord} carries a step size $\alpha^k$, and convergence depends on how that step size is chosen; Diekerhof, for instance, uses a diminishing rule, $\alpha^k=\delta/\sqrt{k}$, which is guaranteed to converge as $k\to\infty$ \cite{Diekerhof2018}. In ADMM, by contrast, a convex problem converges for any $\rho>0$, so $\rho$ affects only the rate: a large value brings the copies into agreement quickly while slowing the multipliers, and a small value does the reverse \cite{Diekerhof2018}. Neither parameter is free in practice, since a poor value slows the method either way, but only the step size enters the condition under which DD is claimed to converge.

The three studies that report timings do not agree on which approach is faster. Lilla et al.\ solve a network of ten prosumers over one day: the decomposed model required about 60\,s in 13 iterations and 80\,s in 8 iterations for its two stages, against about 1\,s per stage for the centralized model, and the authors state that the purpose is not computational performance but the reduction of shared information \cite{Lilla2020}. Diekerhof reports both signs within one thesis: at 10\,000 buildings solved for one scenario, the distributed computation took about 80\,s, almost three times faster than the centralized benchmark, but the order reverses when the same problem covers 27 scenarios and is still solved once \cite{Diekerhof2018}. Olivella-Rosell et al., with 100 sites, report the decomposed solve finishing in under five minutes where the centralized one exhausted 16\,GB of memory and needed about an hour \cite{OlivellaRosell2020}. The three differ in participant count, horizon, treatment of uncertainty, and hardware, so each figure is meaningful only within the study that measured it.


\section{Conclusion}

The comparison yields a choice rule rather than a ranking. Below about fifty prosumer sites the centralized solve is the cheaper option, and above that point the decomposition takes over, reaching a factor of five to twelve at a hundred sites \cite{OlivellaRosell2020}; at ten prosumers the gap runs the other way by two orders of magnitude, about a second for the centralized model against 60\,s and 80\,s for the decomposed one, which solves two problems in sequence \cite{Lilla2020}. The advantage is not only in size: Diekerhof obtains a reversal once the same problem covers 27 scenarios, since each iteration then carries a larger local subproblem \cite{Diekerhof2018}.

That size rule assumes that the aggregator may hold the participants' cost functions and device constraints at all. Where it may not, the question is settled before scale is considered: a decomposition discloses the schedule of each participant and a price for the coupling, not the models behind them \cite{Molzahn2017}, and that restriction is paid for in time. Lilla et al.\ spend precisely this difference on it, taking 60\,s and 80\,s at ten prosumers where the centralized solve takes about a second, in order to reduce the information shared \cite{Lilla2020}. What the aggregator cannot do is check the claim: without the constraints it is not given, it has to take the participants' word that the schedule they return is executable, and that reliance is the price of keeping the models local.

Between the two decompositions the choice rests on what each can promise, not on speed. In the formulation used here, ADMM keeps every iterate on the coupling constraint and, for a convex problem, converges for any positive penalty parameter, so that parameter governs only the rate; in DD the step size enters the convergence condition itself, and the iterates satisfy the coupling only at convergence.

Neither promise survives binary device decisions. Where charging and discharging are mutually exclusive, no decomposition carries an optimum guarantee, and each remedy---convex relaxation, an exact local mixed-integer solve, or heuristic use---pays in feasibility, optimality, or attainable flexibility, with the price reported in no consistent form. What the reviewed studies also leave unmeasured is whether a returned schedule satisfies the discrete constraints it was meant to satisfy. Until that is checked, applicability must be judged per setting, and the rule above holds only within the premise under which the guarantee itself was stated.

% When you add references, uncomment this block and replace the example
% with actual references. Cite entries in the text using \cite{key}.
 \begin{thebibliography}{00}

\bibitem{Diekerhof2018}
M. Diekerhof, ``Distributed optimization for the exploitation
of multi-energy flexibility under uncertainty in city districts,''
Ph.D. dissertation, RWTH Aachen University, Aachen, Germany, 2018.

\bibitem{OlivellaRosell2020}
P. Olivella-Rosell et al., ``Centralised and distributed
optimization for aggregated flexibility services provision,''
IEEE \emph{Trans. Smart Grid}, vol. 11, no. 4, 2020,
doi: 10.1109/TSG.2019.2962269.

\bibitem{Molzahn2017}
D. K. Molzahn et al., ``A survey of distributed optimization
and control algorithms for electric power systems,''
IEEE \emph{Trans. Smart Grid}, vol. 8, no. 6, 2017,
doi: 10.1109/TSG.2017.2720471.

\bibitem{DallAnese2017}
E. Dall'Anese, P. Mancarella, and A. Monti,
``Unlocking flexibility: Integrated optimization and control
of multienergy systems,''
\emph{IEEE Power Energy Mag.}, pp. 43--52, Jan./Feb. 2017,
doi: \url{10.1109/MPE.2016.2625218}.

\bibitem{Boyd2011}
S. Boyd, N. Parikh, E. Chu, B. Peleato, and J. Eckstein,
`Distributed optimization and statistical learning via the alternating
direction method of multipliers,'' \emph{Found. Trends Mach. Learn.},
vol. 3, no. 1, pp. 1--122, 2011, doi: 10.1561/2200000016.

\bibitem{Lilla2020}
S. Lilla, C. Orozco, A. Borghetti, F. Napolitano, and F. Tossani,
``Day-ahead scheduling of a local energy community: An alternating
direction method of multipliers approach,''
IEEE \emph{Trans. Power Syst.}, vol. 35, no. 2, pp. 1132--1142, Mar. 2020,
doi: 10.1109/TPWRS.2019.2944541.

 \end{thebibliography}

\end{document}
